\originalsection{熵、动量通量与能量通量}
密度和速度给四个方程, 可压缩流体还有压强这个未知量. 理想流体无黏性、无热传导时, 每个流体团的比熵$s$满足$Ds/Dt=0$, 以状态方程$p=p(\rho,s)$闭合系统. 这里$s$是每单位质量的熵, 不把整团总熵与随质量标签不变的比熵混用.

连续性与Euler方程合并后, 动量第$i$分量满足
\[
\partial_t(\rho v_i)+\partial_j\Pi_{ij}=\rho f_i,
\qquad \Pi_{ij}=\rho v_iv_j+p\delta_{ij}.
\]
证明是把$\partial_t(\rho v_i)$按乘积展开, 用连续性消去$\partial_t\rho$, 再把$\rho Dv_i/Dt$替成$-\partial_ip+\rho f_i$. $\Pi_{ij}$表示第$i$方向动量通过法向第$j$方向的面的通量, 所以动量密度是向量而通量是张量. 其中$\rho v_iv_j$是流体搬运动量, $p\delta_{ij}$是压力传递动量. 它与只描述面力的Cauchy应力$-p\delta_{ij}$对象不同, 符号差异由其定义决定.

为写能量守恒, 引入比内能$e(\rho,s)$, 热力学微分为$\dd e=T\dd s+(p/\rho^2)\dd\rho$. 沿理想流体团$Ds/Dt=0$, $D\rho/Dt=-\rho\nabla\cdot\vect v$, 故$\rho De/Dt=-p\nabla\cdot\vect v$. 连续性将它变成
\[
\partial_t(\rho e)+\nabla\cdot(\rho e\vect v)
=-p\nabla\cdot\vect v.
\]
Euler方程点乘$\vect v$给动能变化中的$-\vect v\cdot\nabla p$. 两个压力项相加为$-\nabla\cdot(p\vect v)$: 压力既改变体团内能, 也通过体团边界作功. 若体力为$\vect f=-\nabla\Phi$, 总能量密度$w=\rho(v^2/2+e+\Phi)$满足
\[
\partial_tw+\nabla\cdot\bigl[(w+p)\vect v\bigr]
=\rho\partial_t\Phi.
\]
固定外势时右端为零, 积分后得到能量变化等于穿过边界的能量通量. $p/\rho$并非通常意义的比内能; 能量通量中的$e+p/\rho$称比焓, 这一项体现推动流体进入或离开区域所需的功.

若流动等熵且$p$可写成$\rho$的单值函数, 定义$h(\rho)=e+p/\rho$, 则$\dd h=\dd p/\rho$. 定常Euler方程沿流线积分给$v^2/2+h+\Phi=B$, 每条流线有自己的$B$. 常密度不可压缩模型中用$h=p/\rho+$常数, 重得前面的Bernoulli式. 如果有热交换、黏性损失或不同流线熵值, 就须按对应能量方程重新分析, 不能任意两点直接相减.

\originalsection{涡量与流体团的局部转动}
涡量为$\vect\Omega=\nabla\times\vect v$. Stokes积分定理给$\int_S\vect\Omega\cdot\vect n\dd A=\oint_{\partial S}\vect v\cdot\dd\vect l$, 将局部场与环量联系起来. 在一点展开速度梯度, 对称部分改变微小流体团的形状, 反对称部分产生刚体式转动, 局部角速度为$\vect\Omega/2$. 因此剪切流即使没有整体绕轴旋转, 也可有非零涡量.

取无旋体力$\vect f=-\nabla\Phi$. 对Euler方程取旋度, 再展开$\nabla\times(\vect v\times\vect\Omega)$, 可得
\[
\frac{D\vect\Omega}{Dt}
=(\vect\Omega\cdot\nabla)\vect v
-\vect\Omega(\nabla\cdot\vect v)
+\frac{\nabla\rho\times\nabla p}{\rho^2}.
\]
最后一项来自$-\nabla\times(\nabla p/\rho)$, 在等熵正压流体$p=p(\rho)$中为零. 常密度不可压缩理想流体仍有$D\vect\Omega/Dt=(\vect\Omega\cdot\nabla)\vect v$, 即涡线拉伸可以改变涡量. 原稿第123--124页把物质导数与旋度直接交换, 因而称涡量矢量沿流线恒定; 这一论证须修正. 正确的可用结论是: 对受无旋体力的平滑正压理想流体, 初始涡量为零时, 上述齐次输运方程保持它为零. 有黏性且$\rho,\nu$常数时另加$\nu\nabla^2\vect\Omega$, 会使涡量扩散.

\originalsection{静水与定常流的实际应用}
静水$v=0$给$\nabla p=-\rho\nabla\Phi$. 在重力$\Phi=gz$及常密度条件下$p(z)=p(0)-\rho gz$. 对浸没物体积分压力面力, 用散度定理得$\vect F_b=-\int_{V_{\rm displaced}}\nabla p\dd V=\rho gV_{\rm displaced}\hat{\vect z}$, 即Archimedes浮力. 若以等温理想气体$p=\rho R'T$为近似, 则$\dd p/\dd z=-gp/(R'T)$, $p(z)=p(0)e^{-gz/(R'T)}$. 这是所给状态方程与温度模型下的结果, 不应把任何温度梯度都称作必然不稳定.

定常流在各点不随时间变化, 不表示每个流体团都不加速. 变截面水管若近似不可压缩且截面速度均匀, 连续性给$A_1v_1=A_2v_2$, 同高度Bernoulli式给$p_1-p_2=\rho(v_2^2-v_1^2)/2$. 窄处速度较高、静压较低. 大水箱小孔出流时自由液面速度可忽略, 两点同为大气压, 故$v_{\rm jet}=\sqrt{2g(z_{\mathrm{surface}}-z_{\mathrm{hole}})}$. 忽略黏性和孔口收缩才有此理想值.

流线是在同一时刻处处沿$\vect v$的曲线; 迹线是同一个体团满足$\dot{\vect x}=\vect v(\vect x,t)$的运动轨迹. 定常时可消去$t$, 使二者重合; 非定常时一般不同. 驻点$v=0$处不能用$\dd y/\dd x=v_y/v_x$直接作除法. 在撞向驻点的流线上, $p_{\mathrm{stag}}-p_\infty=\rho v_\infty^2/2$. Pitot管比较驻点压与静压, 若压差以密度$\rho_l$的液柱在密度$\rho_a$的气流中测得, 则$\rho_av_\infty^2/2=(\rho_l-\rho_a)g\Delta h$. 液柱是压差读数, 并非直接等于速度.
